\section{NumPy (Iris — session 8)}

\noindent
In session 6 and session 7, we focused on \textbf{classes} and \textbf{inheritance} to build reusable classifiers for the Iris classification task.
This session, we introduce \textbf{NumPy}, the most important Python library for fast numerical computing. NumPy is designed to work efficiently with large arrays of numbers, which is exactly what we need for machine learning.
Recall from session 7 that our classifier interface was \texttt{fit()}, \texttt{predict\_one()}, \texttt{predict\_all()}, and \texttt{evaluate()}.
In this session, we keep that interface but implement NumPy-based data operations underneath it.

\noindent
In this tutorial, you will:
\begin{itemize}\setlength{\itemsep}{0pt}
\item store Iris measurements in NumPy arrays,
\item inspect array properties (\texttt{shape}, \texttt{dtype}),
\item compute a \textbf{centroid} (mean feature vector) for each Iris species,
\item practice importing your own code across multiple files (clean code organization),
\item optionally use an LLM (e.g., ChatGPT) to explain NumPy-heavy code sections.
\end{itemize}

\noindent
\textbf{Deliverable:} \texttt{session8\_iris\_template.py}

\subsection{What You Will Learn}
By the end of this tutorial, you will be able to:
\begin{itemize}\setlength{\itemsep}{0pt}
\item create NumPy arrays and understand their advantages over Python lists for numerical tasks,
\item use indexing and slicing on arrays,
\item compute means with \texttt{np.mean(..., axis=\dots)},
\item use boolean masks to select rows in a dataset,
\item compute Iris class centroids and use them for classification intuition,
\item import custom functions/classes from one Python file to another.
\end{itemize}

\subsection{Tools and Setup}
\noindent
You need:
\begin{itemize}\setlength{\itemsep}{0pt}
\item Python 3,
\item NumPy installed (\texttt{pip install numpy}),
\item an \texttt{iris.csv} file.
\end{itemize}

\noindent
\textbf{Assumption about \texttt{iris.csv}:} it contains 5 columns:
\begin{center}
    \texttt{sepal\_length,sepal\_width,petal\_length,petal\_width,species}
\end{center}

\noindent
If your file has a header row, we will handle it.

\newpage
\subsection{Activity 1: Load Iris into NumPy arrays (\texttt{X} and \texttt{y})}

\noindent
Create a new file called \texttt{session8\_iris\_template.py}. Start with:

\begin{minted}[fontsize=\small, linenos, breaklines=true]{python}
import numpy as np

def load_iris_csv(filepath="iris.csv"):
    """
    Loads iris.csv into:
      X: (n_samples, 4) float array
      y: (n_samples,) string array
    """
    # Use genfromtxt for mixed types (numbers + strings)
    data = np.genfromtxt(
        filepath,
        delimiter=",",
        dtype=None,          # allow mixed dtype
        encoding="utf-8",
        skip_header=1        # assumes there is a header row
    )

    # If your CSV has no header, change skip_header=0
    # data will be a structured array with 5 columns
    X = np.column_stack([data[i] for i in range(4)]).astype(float)
    y = np.array(data[4], dtype=str)

    return X, y

X, y = load_iris_csv("iris.csv")
print("X shape:", X.shape)
print("X dtype:", X.dtype)
print("y shape:", y.shape)
print("y dtype:", y.dtype)
\end{minted}

\subsubsection{Your Task}
\begin{enumerate}
    \item Run the script.
    \item Confirm you see something like \texttt{X shape: (150, 4)} (if using the standard Iris dataset).
    \item Change \texttt{skip\_header} if your CSV format differs.
\end{enumerate}

\newpage
\subsection{Activity 2: Quick NumPy warm-up (indexing and slicing)}

\noindent
Add the following below your prints:

\begin{minted}[fontsize=\small, linenos, breaklines=true]{python}
# First sample (row 0)
print("First row of X:", X[0])
print("First label y[0]:", y[0])

# First 5 rows
print("First 5 rows of X:\n", X[:5])

# Only petal features (columns 2 and 3)
petal = X[:, 2:4]
print("Petal matrix shape:", petal.shape)

# Sepal length column only (column 0)
sepal_length = X[:, 0]
print("Sepal length shape:", sepal_length.shape)
\end{minted}

\subsubsection{Your Task}
\begin{enumerate}
    \item Print the last row using negative indexing.
    \item Print the first 10 labels from \texttt{y}.
\end{enumerate}

\newpage
\subsection{Activity 3: Compute centroids (mean feature vector) per species}

\noindent
A \textbf{centroid} is the mean of the features for a class. For each species, we want:
\[
    \mu_{\text{class}} = \texttt{mean}(X_{\text{class}}, \text{axis}=0)
\]
This gives a 4D vector:
\texttt{[mean\_sepal\_length, mean\_sepal\_width, mean\_petal\_length, mean\_petal\_width]}.

\subsubsection{Step 1: Find unique species}
\begin{minted}[fontsize=\small, linenos, breaklines=true]{python}
classes = np.unique(y)
print("Classes:", classes)
\end{minted}

\subsubsection{Step 2: Compute centroid with a boolean mask}
\begin{minted}[fontsize=\small, linenos, breaklines=true]{python}
centroids = {}

for c in classes:
    mask = (y == c)          # boolean array, True where class matches
    X_c = X[mask]            # select rows belonging to class c
    mu = np.mean(X_c, axis=0)
    centroids[c] = mu

print("Centroids:")
for c in classes:
    print(c, "->", centroids[c])
\end{minted}

\subsubsection{Your Task}
\begin{enumerate}
    \item Print how many samples each class has using \texttt{X\_c.shape[0]}.
    \item For each class centroid, print the centroid rounded to 2 decimals using \texttt{np.round(mu, 2)}.
\end{enumerate}

\newpage
\subsection{Activity 4: Why \texttt{axis=0} matters}

\noindent
In NumPy, \texttt{axis=0} means: “collapse rows (samples) and compute the mean per column (feature)”.

\begin{minted}[fontsize=\small, linenos, breaklines=true]{python}
example = X[:3]  # first 3 samples
print("example shape:", example.shape)
print("mean axis=0 (per feature):", np.mean(example, axis=0))
print("mean axis=1 (per sample):", np.mean(example, axis=1))
\end{minted}

\subsubsection{Your Task}
\begin{enumerate}
    \item In one sentence (as a comment), explain the difference between \texttt{axis=0} and \texttt{axis=1}.
\end{enumerate}

\newpage
\subsection{Activity 5: Code organization across files (imports)}

\noindent
As projects grow, we avoid putting everything in one file. This session, practice splitting code into \textbf{three files}:

\begin{itemize}\setlength{\itemsep}{0pt}
\item \texttt{session8/iris\_session8\_utils.py} (loading + NumPy preparation)
\item \texttt{session8/iris\_session8\_models.py} (session 7-compatible classifier interface with NumPy internals)
\item \texttt{session8/session8\_iris\_template.py} (main script: centroids + evaluation + printing)
\end{itemize}

\noindent
Recall from session 7 that the classifier interface should stay stable.
So in \texttt{session8/iris\_session8\_models.py}, keep the same method names:
\texttt{fit()}, \texttt{predict\_one()}, \texttt{predict\_all()}, and \texttt{evaluate()}.

\subsubsection{Step 1: Create \texttt{session8/iris\_session8\_utils.py}}
\begin{minted}[fontsize=\small, linenos, breaklines=true]{python}
# session8/iris_session8_utils.py
import numpy as np

def load_iris_csv(filepath="iris.csv", skip_header=1):
    data = np.genfromtxt(
        filepath,
        delimiter=",",
        dtype=None,
        encoding="utf-8",
        skip_header=skip_header
    )
    X = np.column_stack([data[i] for i in range(4)]).astype(float)
    y = np.array(data[4], dtype=str)
    return X, y
\end{minted}

\subsubsection{Step 2: Import into \texttt{session8/session8\_iris\_template.py}}
\begin{minted}[fontsize=\small, linenos, breaklines=true]{python}
# session8/session8_iris_template.py
import numpy as np
from session8.iris_session8_utils import load_iris_csv

def main():
    X, y = load_iris_csv("iris.csv", skip_header=1)
    print("X shape:", X.shape, "| dtype:", X.dtype)
    print("y shape:", y.shape, "| dtype:", y.dtype)

    classes = np.unique(y)
    centroids = {}
    for c in classes:
        X_c = X[y == c]
        centroids[c] = np.mean(X_c, axis=0)

    print("\nCentroids (rounded):")
    for c in classes:
        print(c, "->", np.round(centroids[c], 2))

if __name__ == "__main__":
    main()
\end{minted}

\subsubsection{Your Task}
\begin{enumerate}
    \item Run \texttt{session8/session8\_iris\_template.py} and confirm the import works.
    \item Intentionally rename \texttt{session8/iris\_session8\_utils.py} to something else and observe the import error.
    \item Rename it back correctly.
\end{enumerate}

\noindent
This file/module structure is important because Session 9 will import the same session 8 utilities and model interface for plotting-based analysis.

\newpage
\subsection{Optional Activity: Use an LLM to explain NumPy-heavy code (Iris-only)}

\noindent
Professional developers often paste unfamiliar code into an LLM and ask for a guided explanation.
This is a \textbf{learning tool}, not a replacement for understanding.

\subsubsection{Prompt 1: Explain NumPy usage}
\noindent
Copy your centroid computation code and ask:

\begin{quote}
    Explain every place NumPy is used in this script. For each NumPy function, tell me:
    (1) what it does, (2) what shape the input is, (3) what shape the output is.
\end{quote}

\subsubsection{Prompt 2: Explain boolean masks + axis mean}
\noindent
Copy the loop:

\begin{minted}[fontsize=\small, linenos, breaklines=true]{python}
for c in classes:
    X_c = X[y == c]
    centroids[c] = np.mean(X_c, axis=0)
\end{minted}

\noindent
Then ask:

\begin{quote}
    Explain line-by-line how boolean masking works in NumPy here, and why axis=0 produces a centroid vector.
\end{quote}

\subsubsection{Your Task (short reflection)}
Write 3--5 lines (as comments in your Python file):
\begin{itemize}\setlength{\itemsep}{0pt}
\item One thing the LLM explained clearly.
\item One thing you still find confusing.
\end{itemize}

\newpage
\subsection{Submission Requirement (Graded)}

\noindent
Submit:
\begin{enumerate}
    \item \texttt{session8/session8\_iris\_template.py}
    \item \texttt{session8/iris\_session8\_utils.py}
    \item \texttt{session8/iris\_session8\_models.py}
    \item A screenshot/PDF of your terminal output showing:
    \begin{itemize}\setlength{\itemsep}{0pt}
    \item \texttt{X shape} and \texttt{X dtype}
    \item printed centroids for all classes
    \end{itemize}
\end{enumerate}

\vspace{1cm}
\begin{center}
{\color{red}\Large Reminder: Iris-only context. No other datasets.}
\end{center}
